\documentclass[10pt]{article}

\usepackage[utf8]{inputenc}

\usepackage[T1]{fontenc}

\usepackage{amsmath}

\usepackage{amsfonts}

\usepackage{amssymb}

\usepackage[version=4]{mhchem}

\usepackage{stmaryrd}

\usepackage{graphicx}

\usepackage[export]{adjustbox}

\graphicspath{ {./images/} }

\usepackage{bbold}



\begin{document}

ФЕРМА ІІРНЦИП - вариационный принцип, позволяющий находить лучи, т. е. кривые, вдоль к-рых распространяется волновой процесс. Пусть $x=x(\sigma)$, $x=x_{1}, \ldots, x_{m}, \sigma_{0} \leqslant \sigma \leqslant \sigma_{1},-$ уравнение кривой $l$, соединяющей точки $M_{0}$ и $M_{1}, \quad c=c(x)>0$-скорость распространения волн в точке $x$. Ф.п. утверждает, что $\delta \int d s / c=0$ для луча, соединяющего $M_{0}$ и $M_{1}$. Здесь $\delta$-сімвол вариации, $d s=\sqrt{\Sigma\left(d x_{i}\right)^{2}}-$ дифференциал дуги. Фізіч. смысл $\int_{M_{0}}^{M_{1}} d s / c-$ время движения из $M_{0}$ в $M_{1}$ вдоль $l$ со скоростью $c(x)$. Из Ф.п. следуют классич. законы отражения, преломления, прямолинейность лучей при $c=$ const. Диф̆ракционные лучи, лучи, распространяющиеся от краев әкранов, лучи головных волн тоже можно найти, пользуясь Ф.п. Лучи, определяемые Ф.п., являются характеристиками эйконала уравнения. Интеграл $\int d s / c$ задает риманову метрику частного типа. Лучи - геодезические, соответствующие этой метрике. Ф.п. обобщается на случай скорости, зависящей от направления (анизотропная среда). Лучи в этом случае-геодезические нек-рой финслеровой метрики.



Ф. п. для задачи преломления света был впервые высказан П. Ферма (Р. Fermat, ок. 1660).



$\begin{array}{ll}\text { Iит. см. при ст. Лучевой метод. } & \text { В. М. Бабич. }\end{array}$ ФЕРМА СПИРАЛЬ - плоская трансцендентная кривая, уравнение к-рой в полярных координатах имеет вид



\[

\rho=a \sqrt{\bar{\varphi}} .

\]



Каждому значению $\varphi$ соответствуют два значения $\rho-$ положительное и отрицательное. Ф. с. центрально симметрична относительно полюса, к-рый является точкой перегиба. Относится к так называемым алгебраическим спиралям (см. Cиupaлu).



Ф. с. впервые рассмотрена П. Ферма (P., Fermat, 1636).



Лит. $[1]$ С а в е л о в А. А., Плоские кривые, М., 1960. Д. Д. Соколов.



ФEPMA TEOPEMA, в е л и-



\includegraphics[max width=\textwidth, center]{2023_07_09_d9065873d42e9346f76dg-1}

кая те о рема Фер м а, 3 наменитая теорема Ферма, боль лі а т тео рем а Ф е р м а,- утверждение, что для любого натурального числа $n>2$ уравнение $x^{n}+y^{n}=z^{n}$ (у р а в н өни е Ф е р м а) не имеет решений в целых ненулевых числах $x, y, z$. Она была сформулирована П. Ферма (P. Fermat) примерно в 1630 на полях книги Диофанта «Арифметика» [1] следующим образом: «невозможно разложить ни куб на два куба, ни биквадрат на два биквадрата, и вообще никакую степень, больпую квадрата, на две степени с тем жө показателем». И далее добавил: «я открыл әтому поистине чудесное доказательство, но эти поля для него слишком малы». В бумагах П. Ферма напли доказательство Ф. Т. для $n=4$. Общее доказательство до сих пор (1984) не получено, несмотря на усилия многих математиков (как профессионалов, так и любителей). Нездоровый интерес к доказательству әтой теоремы среди неспециалистов в области математики был в свое время вызван большой международной премией, аннулированной в конце 1-й мировой войны.



Предполагается, что доказательство Ф. т. вообще не существовало.



Для $n=3$ Ф. т. доказал Л. Эйлер (L. Euler, 1770), для $n=5$ - П. Дирихле (P. Dirichlet) и А. Лежандр (A. Legendre, 1825), для $n=7$ - Г. Ламе (G. Lamé, 1839) (см. [2]). Ф. т. достаточно доказать для любого простого показателя $n=p>2$, т. е. достаточно доказать, что уравнение



\[

x p+y^{p}=z p

\]



не имеет решений в целых ненулевых взаимно простых числах $x, y, z$. Можно также считать, что числа $x$ и $y$ взаимно просты с $p$. При доказательстве Ф.т. рассматривают два случая: случай 1, когда $(x y z, p)=1$ и случай 2, когда $p \mid z$. Доказательство 2-го случая Ф.т. более сложно и обычно проводится методом бесконечного спуска. Существенный вклад в доказательство Ф.т. внес Ә. Куммер (Е. Китmer), г-рый создал принципиально новый метод, основанный на разработанной им арифметич. теории кругового поля. Используется тот факт, что в поле $\mathbb{Q}(\zeta), \quad \zeta=e^{2 \pi i / p}$, лөвая часть уравнөния (1) разлагается на линейные множители $x^{p}+y^{p}=\prod_{i=0}^{p-1}\left(x+y \zeta^{i}\right)$, к-рые являются $p$-ми степенями идеальньх чисел поля $\mathbb{Q}(\zeta)$ в случае 1 и отличаются от $p-\mathbf{x}$ степеней на множитель $(1-\xi), i>0$ в случае 2. Если $p$ делит числители Бернулии чисел $B_{2 n}(n=1,2, \ldots,(p-3) / 2)$, то по критерию регулярности $p$ не делит число $h$ классов идеалов поля $\mathbb{Q}(\xi)$ и эти идеальные числа-главные. В этом случае Э. Куммер [3] доказал Ф.т. Не известно бесконечно или конечно число регулярных чисел $p$ (по теореме Иенсена число иррегулярных простых чисел бесконетно [4]). Э. Куммер [5] доказал Ф.т. для нек-рых классов иррегулярных простых чисел и тем самым установил ее справедливость для всех $p<100$. В слутае 1 он показал, что из (1) следует выполнимость сравнений



\[

\begin{gathered}

B_{n}\left[\frac{d^{p-n}}{d v^{p-n}} \ln \left(x+e^{v} y\right)\right]_{v=0} \equiv 0(\bmod p), \\

n=2,4, \ldots, p-3

\end{gathered}

\]



справедливых при любой перестановке $x, y,-z . \mathrm{OT}_{-}$ сюда он получил, что если в случае 1 уравнение (1) разрешимо, то для $n=3,5$



\[

B_{p-n} \equiv 0(\bmod p) .

\]



В случае 2 Ә. Куммер доказал Ф.т. при следующих условиях: 1) $p \mid h_{1}, p^{2} \nmid h_{1}$, где $h_{1}$-шервый множитель числа классов идеалов поля $\mathbb{Q}(\zeta)$ (әто равносильно требованию, что числитель только одного числа $B_{2 n}$, где $2 n=2,4, \ldots, p-3$, делится на $p)$; 2) $B_{2 n p} \neq$ $\left.\neq 0\left(\bmod p^{3}\right) ; 3\right)$ найдется идеал, по модулю к-рого единица



\[

E_{n}=\sum_{i=1}^{(p-3) / 2} e_{i}^{g^{-2 i n}}

\]



не сравнима с $p$-й степенью целого числа из $\mathbb{Q}(\zeta)$, где $g$-первообразный корень по модулю $p$, a



\[

e_{i}=\left(\zeta^{g^{i+1} / 2}-\zeta^{-g^{i+1} / 2}\right) /\left(\zeta^{i / 2}-\zeta^{-g^{i} / 2}\right) .

\]



Метод Куммера получил широкое развитие во многих работах по Ф.т. (см. [6], [7]). Из (1) в случае 1 установлена выполнимость с равнений (2) для $n=7,9$, 11, 13, 15, 17, 19. М. Краснер [8] при тех же условиях показал, что найдется такое число $p_{0}$, что при $p>p_{0}$ сравнение (2) верно для всех чисел $n=2 k+1$, где $1 \leqslant k \leqslant[\sqrt[3]{\ln p}]$



Х. Брюкнер [9] доказал, что число чисел $B_{n}, n=$ $=2,4, \ldots, p-3$, qшслители к-рых делятся на $p$ больше, чем $\sqrt{p}-2$. Пусть $p^{k} \mid h_{1}, p^{k+1} \nmid h_{1}$. П. Реморов [10] показал, что существуют такие постоянные $N_{k}$ и $M_{k}$, $N_{k}<M_{k}$, что для всех $p<N_{k}, \quad p>M_{k}$, случай 1 Ф.т. справедлив. М. Айхлер [11] установил, что случай 1 Ф.т. верен при $H \leqslant[\sqrt{p}]-1$, где $H$-индекс иррегулярности поля $\mathbb{Q}(\zeta), p^{H} \| h$. Г. Вандивер [12] доказал случай 1 при $p \nmid h_{2}$, где $h_{2}-$ второй множитель числа классов идеалов поля $\mathbb{Q}(\zeta)$. Интересные резуль-





\end{document}